\section{Reciprocal Rank Fusion outperforms Condorcet and individual Rank Learning Methods}
\paperhead{Reciprocal Rank Fusion outperforms Condorcet and individual Rank Learning Methods}{G. V. Cormack, C. L. A. Clarke, Stefan B\"{u}ttcher}{ACM SIGIR 2009, pages 758 to 759}{https://doi.org/10.1145/1571941.1572114}{OpenAlex title match 1.00, cited\_by=642}

\subsection{What the paper does}

The paper attacks metaranking: given several IR systems that each return a ranked list for one query, produce a combined list that beats all of them. Supervised learning to rank was the fashionable answer but needs labelled data. The established unsupervised standards were Condorcet Fuse, which sorts documents by pairwise majority vote across the input rankings, and CombMNZ, which sums each system's own score for a document and multiplies by the number of systems ranking it above a cutoff $c$, so it depends on scores whose scales are arbitrary and mutually incomparable.

The authors devised Reciprocal Rank Fusion as a throwaway baseline for that comparison and found it beat the methods it was meant to calibrate. A document's fused score is a sum of reciprocals of its ranks, using rank alone and never the producing system's score. The additive constant $k$ damps how much one system's top positions can dominate, while reciprocal decay keeps deep ranks contributing something, which an exponential decay would not.

\subsection{How the method works}

Given a document set $D$ and a set $R$ of rankings, each a permutation on $1..|D|$, the score is
\[ \mathrm{RRFscore}(d \in D) = \sum_{r \in R} \frac{1}{k + r(d)}, \qquad k = 60 . \]
The value $k=60$ was fixed during a pilot investigation and explicitly not altered during subsequent validation. Three experiment families follow. First, four pilot runs, each fusing 30 configurations of the Wumpus search engine over four different TREC collections; the $k$ sweep is reported for TREC topics 351 to 400. Second, four sets of runs actually submitted by TREC participants (TREC 3, TREC 5 and TREC 9 ad hoc, plus the TREC 2004 Robust track), chosen because prior metaranking work used them. Third, a meta-learner built by applying RRF to the baseline learning-to-rank results shipped with LETOR 3. Significance comes from a sign test across runs, plus confidence intervals and $p$ values on LETOR.

Two operational properties matter as much as the formula. RRF combines ranks without regard to the arbitrary scores returned by particular ranking methods, so no score normalisation is required; and it needs no voting algorithm or global information, since ranks may be computed and summed one system at a time rather than held in memory together.

\subsection{What it reports}

In the pilot, MAP rises from .2072 at $k=0$ to .2145 at $k=60$, peaks at .2147 at $k=80$ and falls to .2098 at $k=500$, which the authors read as $k=60$ being near optimal but the choice not being critical; the best individual model there scored .2016, Condorcet .2074, CombMNZ .2039. On submitted TREC runs RRF scores .3686 on Robust, .4350 on TREC 3, .3394 on TREC 5 and .2830 on TREC 9, beating Condorcet in all four cases and CombMNZ in all but TREC 3, where CombMNZ reaches .4381. On TREC 9 the best individual run (.3519) used a human in the loop; RRF beats the best automated run at .2801. Overall the authors claim RRF outperforms Condorcet, CombMNZ and the best system by 4\% to 5\% on average, with sign tests giving $p \approx 0.008$ against Condorcet and $p \approx .04$ against CombMNZ. On LETOR 3, over 583,850 document-query pairs, RRF reaches MAP 0.6051 against ListNet 0.5846, AdaRank-MAP 0.5778, RankSVM 0.5737 and RankBoost 0.5622, bettering every individual ranker at $p < .003$; CombMNZ edges it at 0.6107, though not significantly ($p \approx .2$). Only MAP is tabulated.

\subsection{How Mytheus uses this}
\begin{takeaway}
This paper is the sole origin of the $1/(k+\mathrm{rank})$ score and the $k=60$ constant frozen into the Mytheus fuse-and-rerank stage, so it must be cited there verbatim. The inheritance is unsafe in one specific way. Cormack fuses 30 homogeneous ad hoc systems, all producing relevance orderings over one corpus; Mytheus fuses exactly three heterogeneous arms, and the AGE arm's position is a traversal ordering over RELATES\_TO and FOLLOWS\_FROM lineage with no calibrated relevance semantics at all. With three arms rather than thirty a single arm swings the fused list far harder, so the flatness of MAP across $k$ is evidence about a 30 system ensemble, not a 3 arm one. Build the weighted variant $\sum_a w_a / (k + r_a(d))$ and run the ablation the paper itself ran: sweep $k$ over $\{0, 10, 20, \ldots, 100, 500\}$ on the golden set in CI, measure recall@50 of the fused top-50 and NDCG@10 after the BGE-reranker-base cross-encoder, and repeat with each arm dropped. That is the cheapest defensible experiment in the retrieve stage and the most likely reviewer challenge. RRF survives the latency budget trivially: integer arithmetic over three short lists, negligible against the warm p95 under 500 ms target or the 1.5 s ContextGapAgent round trip, and its one system at a time accumulation maps directly onto three arms landing asynchronously, so a slow arm can be dropped at deadline without recomputing. Two things it cannot give: it is query and user independent, so per-participant personalisation must stay upstream in the expand stage variants or downstream in the reranker and cannot be smuggled into fusion; and since it discards scores, the fused value is not a calibrated relevance probability and must never feed the 0.7 and 0.4 confidence gates without separate calibration.
\end{takeaway}

\subsection{Limitations and what it leaves open}
\begin{itemize}[leftmargin=*, itemsep=1pt, topsep=2pt]
  \item $k=60$ is tuned once, on one pilot ensemble of 30 Wumpus configurations over TREC topics 351 to 400, and deliberately never retuned; the paper establishes insensitivity to $k$ only for that large homogeneous ensemble and never for a small fusion of two or three arms, which is the regime every modern hybrid retriever actually operates in.
  \item Every fused list is a relevance ordering produced by an ad hoc document retrieval system, so the paper never tests an arm whose rank means something else (graph traversal depth, lexical match count) nor an arm that is simply uninformative for a given query, and it offers no diagnostic for detecting such an arm.
  \item All systems contribute with equal weight and are always included: there is no per-arm weighting, no per-query arm selection or dropping, and no study of list truncation depth, even though CombMNZ's explicit cutoff $c$ shows the authors were aware the depth choice matters.
  \item Fusion is unsupervised and produces one identical ranked list for every user, with no personalisation signal, no per-tenant or per-collection partitioning, and no mechanism for turning the fused score into a calibrated relevance value usable by a downstream threshold.
  \item Evaluation is confined to static, pooled TREC and LETOR collections with no latency, memory or cost measurement, no incremental or streaming setting in which the candidate set changes mid session, and only MAP tabulated, with P@k, R-precision and NDCG asserted comparable but never reported.
\end{itemize}
